% Lecture 2 — Groups and fields
% Compiled to web slides by tools/tex2slides.mjs

\lecture{Groups and fields}{Linear Algebra · Warsaw University of Technology}{Lecture 2 · Tomasz Brengos}

\section{Motivation}

\begin{frame}{An impossible-sounding problem}
Every time your browser shows a padlock, this just happened:

\pause

Two computers that have \emph{never met} agree on a secret key —
while every message between them travels through wires and routers
that anyone can read.

\pause

\begin{warning}
Whisper a secret across a crowded room, in plain voice, so that everyone
hears every word — yet only the two of you learn the secret. That is the
requirement.
\end{warning}

\pause

It is solved billions of times a day. The solution is not locks or
wires — it is \emph{algebra}: arithmetic in a finite group.
\note{The padlock hook. Emphasise: no shared secret beforehand, public channel only.}
\end{frame}

\begin{frame}{Diffie--Hellman, with real numbers (1976)}
Everything in \textbf{bold} is public. Alice keeps $a = 6$ secret,
Bob keeps $b = 15$ secret.

\begin{steps}
\step[public agreement, everyone sees it]{\textbf{p = 23}, \qquad \textbf{g = 5}}
\step[Alice sends her mix — the exponent stays hidden]{A = g^a \bmod p = 5^6 \bmod 23 = \textbf{8}}
\step[Bob sends his mix]{B = g^b \bmod p = 5^{15} \bmod 23 = \textbf{19}}
\step[each raises the other's mix to their own secret]{B^a \equiv g^{ba} = g^{ab} \equiv A^b \equiv 2 \pmod{23}}
\end{steps}

\pause

Both arrive at the shared secret $2$. The eavesdropper holds $23, 5, 8, 19$ —
and recovering $a$ from $8 = 5^a \bmod 23$ is the \emph{discrete logarithm
problem}: easy to state, and (for a 600-digit $p$) believed to be hopelessly
hard. The entire padlock rests on that asymmetry.
\end{frame}

\begin{frame}{The algebra your day already used}
\begin{itemize}
\item \textbf{Key exchange and RSA} — arithmetic modulo a prime: the finite
groups and fields of this lecture.
\pause
\item \textbf{QR codes, Wi-Fi, scratched CDs} — Reed--Solomon error correction
and the AES cipher compute in a field with $256$ elements.
\pause
\item \textbf{Every single bit} — $\{0,1\}$ with XOR and AND is the smallest
field of all. Hashing, checksums and cryptanalysis are linear algebra over it.
\end{itemize}
\end{frame}

\begin{frame}{Programming against an interface}
An \emph{interface} declares which operations a type must support and which
laws they must obey — and code written against the interface runs on
\emph{every} implementation.

\pause

This lecture does exactly that to arithmetic:

\begin{itemize}
\item a \emph{group} — one operation, three laws;
\item a \emph{field} — two operations, tied by distributivity.
\end{itemize}

\pause

\begin{remark}
Every theorem proved from the axioms alone is proved \emph{once} — for
$\R$, for $\C$, for residues modulo a prime, for invertible matrices.
And fields are precisely the \emph{scalars} over which the whole course
is built: vector spaces, matrices, determinants, eigenvalues.
\end{remark}
\end{frame}

\section{Groups}

\begin{frame}{Definition}
\begin{definition}
A \emph{group} is a set $G$ with a binary operation $*\colon G\times G\to G$
such that:
\begin{itemize}
\item $(a*b)*c = a*(b*c)$ \hfill (associativity)
\item there is $e\in G$ with $a*e = e*a = a$ for all $a$ \hfill (identity)
\item every $a\in G$ has $a^{-1}\in G$ with $a*a^{-1} = a^{-1}*a = e$ \hfill (inverses)
\end{itemize}
If moreover $a*b = b*a$ for all $a,b$, the group is \emph{abelian}.
\end{definition}

\pause

\begin{remark}
Additive notation: operation $+$, identity $0$, inverse $-a$.
Multiplicative notation: operation $\cdot$, identity $1$, inverse $a^{-1}$.
\end{remark}
\end{frame}

\begin{frame}{Modular arithmetic}
\begin{definition}
For a positive integer $n$, $\ a \bmod n$ is the remainder of dividing $a$
by $n$: the unique $r \in \{0, 1, \ldots, n-1\}$ with $a = qn + r$.
On $\Z_n = \{0, 1, \ldots, n-1\}$ we define
\[
a +_n b = (a+b) \bmod n, \qquad a \cdot_n b = (a \cdot b) \bmod n.
\]
\end{definition}

\pause

\begin{example}
$17 \bmod 4 = 1$, \quad $4 \bmod 17 = 4$, \quad $(-2) \bmod 5 = 3$
(the remainder is always in $\{0,\ldots,n-1\}$).
In $\Z_6$: \ $4 +_6 5 = 3$ and $4 \cdot_6 5 = 2$.
\end{example}

\pause

This is clock arithmetic — hour $9$ plus $5$ hours is hour $2$ — and it is
the arithmetic Diffie--Hellman ran on.
\end{frame}

\begin{frame}{A zoo of groups}
\begin{itemize}
\item $(\Z, +)$, $(\Q, +)$, $(\R, +)$, $(\C, +)$ — abelian; identity $0$,
inverse $-a$.
\pause
\item $(\Q\setminus\{0\}, \cdot)$, $(\R\setminus\{0\}, \cdot)$,
$(\C\setminus\{0\}, \cdot)$ — abelian; identity $1$, inverse $\tfrac1a$.
\pause
\item $(\Z, \cdot)$ is \emph{not} a group: apart from $\pm 1$, no integer has
an integer reciprocal.
\pause
\item $(\Z_n, +_n)$ — a finite abelian group: the inverse of $a$ is $n - a$.
\pause
\item Invertible $2\times 2$ matrices under multiplication (Lecture 4) — a
group that is \emph{not} abelian.
\end{itemize}
\end{frame}

\begin{frame}{Finite multiplicative groups}
For a prime $p$, the non-zero residues $\Z_p \setminus \{0\}$ form a group
under $\cdot_p$ — the group where Diffie--Hellman lives.

\pause

Inverses let us solve equations exactly as in $\R$. In $\Z_7$, solve
$2 \cdot_7 x = 5$:

\begin{steps}
\step[find the inverse of 2]{2 \cdot_7 4 = 8 \bmod 7 = 1, \qquad \text{so } 2^{-1} = 4}
\step[multiply both sides by it]{x = 4 \cdot_7 5 = 20 \bmod 7 = 6}
\step[check]{2 \cdot_7 6 = 12 \bmod 7 = 5 \checkmark}
\end{steps}
\end{frame}

\begin{frame}{Roots of unity revisited}
\begin{columns}
\column
The $n$-th roots of unity from Lecture 1,
\[
U_n = \left\{ \cos\tfrac{2k\pi}{n} + i \sin\tfrac{2k\pi}{n} \right\},
\]
form an abelian group under multiplication: arguments add, so the product of
two roots is again a root.

Identity: $1$. Inverse: the conjugate.

\pause

Multiplying the roots adds the exponents $k$ modulo $n$ — the group $U_n$
\emph{is} $(\Z_n, +_n)$ wearing a geometric costume.
\column
\widget{RootsOfUnity}{n=6}
\end{columns}
\end{frame}

\begin{frame}{First theorems from the axioms alone}
\begin{fact}
In any group the identity is unique, and each element has exactly one inverse.
\end{fact}

\begin{steps}
\step[two identities e and e — prime]{e = e * e' = e'}
\step[two inverses b and b — prime — of a]{b = b * e = b * (a * b')}
\step[associativity finishes it]{= (b * a) * b' = e * b' = b'}
\end{steps}

\pause

\begin{fact}
Cancellation: $a * b = a * c$ implies $b = c$ (multiply by $a^{-1}$ on the
left). These proofs used nothing but the three axioms — so they hold in
\emph{every} group in the zoo, matrices included.
\end{fact}
\end{frame}

\section{Fields}

\begin{frame}{Definition}
\begin{definition}
A \emph{field} is a set $\K$ with two operations $+$ and $\cdot$ such that:
\begin{itemize}
\item $(\K, +)$ is an abelian group with identity $0$;
\item $(\K \setminus \{0\}, \cdot)$ is an abelian group with identity $1$;
\item multiplication distributes over addition: $(a+b)\cdot c = a\cdot c + b\cdot c$.
\end{itemize}
\end{definition}

\pause

Subtraction and division come for free: $a - b = a + (-b)$ and
$\tfrac{a}{b} = a \cdot \tfrac1b$ for $b \neq 0$. All four operations,
with all the usual rules.

\pause

\begin{example}
$\Q$, $\R$ and $\C$ are fields. $\Z$ is \emph{not}: $2$ has no integer
reciprocal.
\end{example}
\end{frame}

\begin{frame}{First consequences}
\begin{fact}
In any field: $a \cdot 0 = 0$, \quad $(-1)\cdot a = -a$, \quad and there are
\emph{no zero divisors}:
\[
a \cdot b = 0 \implies a = 0 \ \text{ or } \ b = 0.
\]
\end{fact}

\pause

\begin{steps}
\step[suppose ab = 0 and a is nonzero]{b = 1 \cdot b = \left( \tfrac1a \cdot a \right) \cdot b}
\step[associativity, then the first fact]{= \tfrac1a \cdot (a \cdot b) = \tfrac1a \cdot 0 = 0}
\end{steps}

\pause

No zero divisors is exactly what makes cancellation — and solving
equations — possible.
\end{frame}

\begin{frame}{The multiplication table decides everything}
\begin{columns}
\column
Is $\Z_n$ a field? The table answers.

An element $a$ is \emph{invertible} exactly when a $1$ appears in its row —
highlighted below.

\textbf{Try it:} set $n = 5$, then $n = 4$. Watch the row of $2$: modulo $4$
it reads $0, 2, 0, 2$ — no $1$, ever. And $2 \cdot_4 2 = 0$: a zero divisor.

Switch to $+$: \emph{every} $\Z_n$ is a group under addition.
\column
\widget{ModTable}{n=5, op=mul}
\end{columns}
\end{frame}

\begin{frame}{When is it a field?}
\begin{theorem}
$\Z_n$ is a field \ if and only if \ $n$ is a prime number.
\end{theorem}

\pause

\begin{steps}
\step[composite n gives zero divisors]{n = a \cdot b, \quad 1 < a, b < n \implies a \cdot_n b = 0}
\step[prime p: the row of a has no zeros and no repeats]{a \cdot_p x = a \cdot_p y \implies p \mid a(x - y) \implies x = y}
\step[n minus 1 distinct non-zero products must fill the row]{\{ a \cdot_p 1, \ldots, a \cdot_p (p-1) \} = \{1, \ldots, p-1\} \ni 1}
\end{steps}

\pause

\begin{remark}
The proof shows an inverse \emph{exists} without finding it. In practice
inverses modulo $p$ are computed by the extended Euclidean algorithm — a
staple of every crypto library.
\end{remark}
\end{frame}

\begin{frame}{The smallest field: a single bit}
\begin{example}
$\Z_2 = \{0, 1\}$ with the tables
\[
\begin{array}{c|cc} + & 0 & 1 \\ \hline 0 & 0 & 1 \\ 1 & 1 & 0 \end{array}
\qquad\qquad
\begin{array}{c|cc} \cdot & 0 & 1 \\ \hline 0 & 0 & 0 \\ 1 & 0 & 1 \end{array}
\]
is a field — the smallest one.
\end{example}

\pause

Read the tables again: addition is the \textbf{XOR} gate, multiplication is
the \textbf{AND} gate. A bit is not a metaphor for a field element —
it \emph{is} one.

\pause

\begin{remark}
Linear algebra over $\Z_2$ powers checksums, hashing, error-correcting codes
and cryptanalysis. AES and QR codes use its big sibling: a field with
$2^8 = 256$ elements, one per byte.
\end{remark}
\end{frame}

\begin{frame}{A field between the rationals and the reals}
\begin{example}
$\Q(\sqrt2) = \{ a + b\sqrt2 \mid a, b \in \Q \}$ is a field.
\end{example}

The only doubtful axiom is the reciprocal — and the trick is the one that
worked for complex numbers:

\begin{steps}
\step[multiply by the conjugate]{\frac{1}{a + b\sqrt2} = \frac{a - b\sqrt2}{(a + b\sqrt2)(a - b\sqrt2)}}
\step[the denominator is rational]{= \frac{a}{a^2 - 2b^2} - \frac{b}{a^2 - 2b^2}\,\sqrt2}
\step[and nonzero, because sqrt 2 is irrational]{a^2 - 2b^2 \neq 0 \ \text{ for } (a,b) \neq (0,0)}
\end{steps}
\end{frame}

\begin{frame}{Summary}
\begin{itemize}
\item A \textbf{group}: one operation, three laws — associativity, identity,
inverses. Proofs from the axioms hold for every implementation at once.
\pause
\item \textbf{Modular arithmetic}: $\Z_n$ is always an additive group;
its non-zero part is a multiplicative group exactly for prime $n$.
\pause
\item A \textbf{field}: two intertwined abelian groups — $\Q$, $\R$, $\C$,
$\Z_p$, $\Q(\sqrt2)$, and the single bit $\Z_2$.
\pause
\item Cryptography, error correction and your Wi-Fi all compute in finite
groups and fields — this is applied mathematics at planetary scale.
\pause
\item \textbf{Next lecture:} vector spaces — where fields become the
\emph{scalars} and linear algebra properly begins.
\end{itemize}
\end{frame}
